% -*- coding: utf-8 -*-
% !TEX program = xelatex
\documentclass[plain,noanswer,medmath]{../../template/jnuexam}
\usepackage{../../template/kysx}

\begin{document}


\examtitle{2014}{1}


\makepart{选择题}{1～8小题，每小题4分，共32分}


\begin{problem}
下列曲线有渐近线的是\pickout{}
\begin{abcd}
\item $y = x + \sin x$．
\item $y = x^2 + \sin x$．
\item $y = x + \sin \frac{1}{x}$．
\item $y = x^2 + \sin \frac{1}{x}$．
\end{abcd}
\end{problem}


\begin{problem}
设函数$f(x)$具有二阶导数，$g(x) = f(0)(1 - x) + f(1)x$，则在区间$[0,1]$上\pickout{}
\begin{abcd}
\item 当$f'(x) \ge 0$时，$f(x) \ge g(x)$．
\item 当$f'(x) \ge 0$时，$f(x) \le g(x)$．
\item 当$f''(x) \ge 0$时，$f(x) \ge g(x)$．
\item 当$f''(x) \ge 0$时，$f(x) \le g(x)$．
\end{abcd}
\end{problem}


\begin{problem}
设$f(x)$是连续函数，则$\int_0^1 \dy\int_{- \sqrt{1 - y^2}}^{1 - y} f(x,y)\dx =$\pickout{}
\begin{abcd}
\item $\int_0^1 \dx\int_0^{x - 1} f(x,y)\dy + \int_{-1}^0 \dx\int_0^{\sqrt{1 - x^2}} f(x,y)\dy$．
\item $\int_0^1 \dx\int_0^{1 - x} f(x,y)\dy + \int_{-1}^0 \dx\int_{- \sqrt{1 - x^2}}^0 f(x,y)\dy$．
\item $\int_0^{\frac{\pi}{2}}\d\theta
       \int_0^{\frac{1}{\cos\theta+\sin\theta}}f(r\cos\theta,r\sin\theta)\dr
       + \int_{\frac{\pi}{2}}^{\pi}\d\theta \int_0^1 f(r\cos\theta,r\sin\theta)\dr$．
\item $\int_0^{\frac{\pi}{2}}\d\theta 
       \int_0^{\frac{1}{\cos\theta+\sin\theta}}f(r\cos\theta,r\sin\theta)r\dr
       + \int_{\frac{\pi}{2}}^{\pi}\d\theta \int_0^1 f(r\cos\theta,r\sin\theta)r\dr$．
\end{abcd}
\end{problem}


\begin{problem}
若$\int_{- \pi}^{\pi}(x - a_1\cos x - b_1\sin x)^2 \dx
= \min\limits_{a,b \in \mathbb{R}}\left\{\int_{-\pi}^{\pi}(x - a\cos x - b\sin x)^2 \dx\right\}$，
\goodbreak 则$a_1\cos x + b_1\sin x =$\pickout{}
\begin{abcd}
\item $2\sin x$．
\item $2\cos x$．
\item $2\pi \sin x$．
\item $2\pi \cos x$．
\end{abcd}
\end{problem}


\begin{problem}
行列式$\begin{vmatrix}
  0 & a & b & 0 \\
  a & 0 & 0 & b \\
  0 & c & d & 0 \\
  c & 0 & 0 & d
\end{vmatrix} =$\pickout{}
\begin{abcd}
\item $(ad - bc)^2$．
\item $- (ad - bc)^2$．
\item $a^2d^2 - b^2c^2$．
\item $b^2c^2 - a^2d^2$．
\end{abcd}
\end{problem}


\begin{problem}
设$a_1,a_2,a_3$均为$3$维向量，则对任意常数$k,l$，向量组$a_1 + ka_3$,$a_2 + la_3$
线性无关是向量组$B = (\alpha_1,\alpha_2,\alpha_3)$线性无关的\pickout{}
\begin{abcd}
\item 必要非充分条件．
\item 充分非必要条件．
\item 充分必要条件．
\item 既非充分也非必要条件．
\end{abcd}
\end{problem}


\begin{problem}
设随机事件$A$与$B$相互独立，且$P(B) = 0.5$，$P(A - B) = 0.3$，则$P(B - A) =$\pickout{}
\begin{abcd}
\item $0.1$．
\item $0.2$．
\item $0.3$．
\item $0.4$．
\end{abcd}
\end{problem}


\begin{problem}
设连续性随机变量$X_1$与$X_2$相互独立，且方差均存在，$X_1$与$X_2$的概率密度分别为$f_1(x)$与$f_2(x)$，
随机变量$Y_1$的概率密度为$f_{Y_1}(y) = \frac{1}{2}[f_1(y) + f_2(y)]$，
随机变量$Y_2 = \frac{1}{2}(X_1 + X_2)$，则\pickout{}
\begin{abcd}
\item $EY_1 > EY_2$，$DY_1 > DY_2$．
\item $EY_1 = EY_2$，$DY_1 = DY_2$．
\item $EY_1 = EY_2$，$DY_1 < DY_2$．
\item $EY_1 = EY_2$，$DY_1 > DY_2$．
\end{abcd}
\end{problem}


\makepart{填空题}{9～14小题，每小题4分，共24分}


\begin{problem}
曲面$z = x^2(1 - \sin y) + y^2(1 - \sin x)$在点$(1,0,1)$处的切平面方程为 \fillin{}．
\end{problem}


\begin{problem}
设$f(x)$是周期为$4$的可导奇函数，且$f'(x)= 2(x - 1)$, $x \in[0,2]$，则$f(7) =$ \fillin{}．
\end{problem}


\begin{problem}
微分方程$xy' + y(\ln x - \ln y) = 0$满足条件$y(1) = \e^3$的解为$y =$ \fillin{}．
\end{problem}


\begin{problem}
设$L$是柱面$x^2 + y^2 = 1$与平面$y + z = 0$的交线，从$z$轴正向往$z$轴负向看去为逆时针方向，
则曲线积分$\oint_L z\dx + y\dz =$ \fillin{}．
\end{problem}


\begin{problem}
设二次型$f(x_1,x_2,x_3) = x_1^2 - x_2^2 + 2ax_1x_3 + 4x_2x_3$的负惯性指数是$1$，
则$a$的取值范围 \fillin{}．
\end{problem}


\begin{problem}
设总体$X$的概率密度为$f\left(x;\theta \right) = \begin{cases}
  \frac{2x}{3\theta^2}, & \theta < x < 2\theta , \\
                     0, & \text{其他},
\end{cases}$其中$\theta$是未知参数，$X_1$, $X_2$, $\cdots$, $X_n$为来自总体$X$的简单样本，
若$c\sum_{i = 1}^n X_i^2$是$\theta^2$的无偏估计，则$c =$ \fillin{}．
\end{problem}


\makepart{解答题}{15～23小题，共94分}


\begin{problem}[points=10]
求极限$\lim_{x \to +\infty}
\frac{\int_1^x \left[t^2\left(\e^{\frac{1}{t}} - 1\right) - t \right]\dt}
     {x^2\ln \left(1 + \frac{1}{x}\right)}$．
\end{problem}


\begin{problem}[points=10]
设函数$y = f(x)$由方程$y^3 + xy^2 + x^2 y + 6 = 0$确定，求$f(x)$的极值．
\end{problem}


\begin{problem}[points=10]
设函数$f(u)$具有二阶连续导数，$z = f\left(\e^x\cos y \right)$满足
$$\frac{\pd^2 z}{\pd x^2} + \frac{\pd^2 z}{\pd y^2} = \left(4z + \e^x\cos y \right) \e^{2x}.$$
若$f(0) = 0,f'(0) = 0$，求$f(u)$的表达式．
\end{problem}


\begin{problem}[points=10]
设$\Sigma$为曲面$z = x^2 + y^2$$(z \le 1)$的上侧，计算曲面积分
$$I = \iint_{\Sigma}(x - 1)^3 \dy\dz + (y - 1)^3 \dz\dx + (z - 1)\dx\dy.$$
\end{problem}


\begin{problem}[points=10]
设数列$\left\{a_n \right\},\left\{b_n \right\}$
满足$0 < a_n < \frac{\pi}{2}$，$0 < b_n < \frac{\pi}{2}$，$\cos a_n - a_n = \cos b_n$，
且级数$\sum_{n = 1}^{\infty}b_n$收敛．
\begin{enumerate*}
\item 证明：$\lim_{n \to \infty} a_n = 0$．
\item 证明：级数$\sum_{n = 1}^{\infty}\frac{a_n}{b_n}$收敛．
\end{enumerate*}
\end{problem}


\begin{problem}[points=11]
设矩阵$A = \begin{pmatrix}
  1 & -2 &  3 & -4 \\
  0 &  1 & -1 & 1 \\
  1 &  2 &  0 & 3
\end{pmatrix}$，$E$为$3$阶单位矩阵．\par
\begin{enumerate*}
\item 求方程组$Ax = 0$的一个基础解系；
\item 求满足$AB = E$的所有矩阵$B$．
\end{enumerate*}
\end{problem}


\begin{problem}[points=11]
证明$n$阶矩阵$\begin{pmatrix}
       1 &      1 & \cdots & 1 \\
       1 &      1 & \cdots & 1 \\
  \vdots & \vdots & \vdots & \vdots \\
       1 &      1 & \cdots & 1
\end{pmatrix}$与$\begin{pmatrix}
       0 & \cdots &      0 & 1 \\
       0 & \cdots &      0 & 2 \\
  \vdots & \vdots & \vdots & \vdots \\
       0 & \cdots &      0 & n
\end{pmatrix}$相似．
\end{problem}


\begin{problem}[points=11]
设随机变量$X$的概率分布为$P\{X = 1\} = P\{X = 2\} = \frac{1}{2}$，
在给定$X = i$的条件下，随机变量$Y$服从均匀分布$U(0,i)$（$i = 1,2$）．\par
\begin{enumerate*}
\item 求$Y$的分布函数$F_Y(y)$；
\item 求$EY$．
\end{enumerate*}
\end{problem}


\begin{problem}[points=11]
设总体$X$的分布函数为$F(x;\theta) = \begin{cases}
  1 - \e^{- \frac{x^2}{\theta}}, & x \ge 0, \\
                              0, & x < 0,
\end{cases}$其中$\theta$是未知参数且大于零．$X_1,X_2,\cdots,X_n$为来自总体$X$的简单随机样本．
\begin{enumerate}
\item 求$E(X)$与$E(X^2)$；
\item 求$\theta$的最大似然估计量$\hat{\theta}_n$；
\item 是否存在实数$a$，使得对任何$\epsilon > 0$，都有
      $\lim_{n \to \infty} P\left\{\left|\hat{\theta}_n - a\right| \ge \epsilon\right\} = 0$？
\end{enumerate}
\end{problem}


\end{document}
